\section{Setup and review of the Oaxaca--Blinder decomposition (OBD)}
\label{sec:counterfactual_OBD}

We first review the index-number problem and the OBD.
We start by considering the case where we have perfect population information about 
all variables of interest. Then we describe how to proceed with real data.


%In this section, we revisit the OBD from a counterfactual perspective. 
%We show how this view naturally leads to index-number problem; i.e., that the choice of reference group in the OBD is arbitrary. 
%We first show how the index-number problem arises at the level of counterfactual mean comparisons, before introducing the standard linear formulation of the OBD.

\input{Sections/02a_counterfactual_explanation.tex}

% \input{Sections/02b_setup}

\subsection{Real-data Oaxaca--Blinder decomposition and drawing conclusions} 
\label{sec:ob_real}

In general, data analysts will have access to real data for each of the two groups. 
\textcolor{red}{We assume they estimate $\mathbb{E}_g[Y|X]$ separately in each group with $\hat{M}_g(X)$, such as ordinary least squares (OLS) or an ML algorithm, and estimate by $\mathbb{E}_{g'}[\mathbb{E}_g[Y|X]]$ by taking the sample average of $\hat{M}_g(X)$ in data from $g'$.}

It is common to draw substantive conclusions based on the signs of the empirical explained and unexplained components, often coupled with
statistical significance \citep{oaxacaRansom1998, jann2008,fortinLemieuxFirpo2011}. For instance, consider the hospital example. 
Suppose now that Hospital $H$ has lower mortality ($\Delta_Y < 0$), and suppose we take $H$ as the reference. 
If the OBD explained component has a negative sign, we might conclude that group $H$ can attribute
its lower mortality rate in part to its risk factors, such as lower blood pressure. In real data, we might additionally require that the observed sign be statistically significant in order to form a conclusion;
without significance, we might
instead hesitate to conclude that risk factors are driving any difference in mortality.
Analogous to the explained component, if the OBD unexplained component has a negative sign, we might attribute group $H$'s lower
mortality in part to the quality of hospital care.
